\section{曲面：第一与第二基本形式}
\label{sec:12-Real-Analysis-Support/11-Curve-and-Surface-Geometry/02}

拿一张纸，上面画好网格和一个三角形，然后把纸卷成圆筒。三角形的三条边长没变，三个角没变，网格的每个小方块面积也没变——纸没被拉伸过一丝一毫。可是这个圆筒摆在桌上，谁都看得出它是弯的。

同一件东西，说它弯也对，说它没弯也对，取决于你站在哪里问。站在纸面上量长度和角度，什么都没发生；站在房间里看它的形状，明显变了。曲线只需要一个数就能说完弯曲，曲面显然不行——这里至少有两种弯曲，需要两套工具分别记账。

\begin{center}
\begin{tikzpicture}[>=stealth,scale=0.9]
% 马鞍面网格（手工二维投影：u 方向上弯，v 方向下弯）
\foreach \s in {-1.6,-0.8,0.8,1.6} {
  \draw[gray!35] plot[domain=-1.8:1.8,samples=40]
    ({\x+0.75*\s},{-0.25*\x+0.30*\x*\x+0.42*\s-0.30*\s*\s});
}
\foreach \t in {-1.8,-0.9,0.9,1.8} {
  \draw[gray!35] plot[domain=-1.6:1.6,samples=40]
    ({\t+0.75*\x},{-0.25*\t+0.30*\t*\t+0.42*\x-0.30*\x*\x});
}
% 过原点的两条法截线
\draw[orange!85!black,line width=1.3pt] plot[domain=-1.8:1.8,samples=50]
  ({\x},{-0.25*\x+0.30*\x*\x});
\draw[blue!70!black,line width=1.3pt] plot[domain=-1.6:1.6,samples=50]
  ({0.75*\x},{0.42*\x-0.30*\x*\x});
% 法向
\draw[->,line width=1pt] (0,0) -- (-0.18,1.35);
\node[left=1pt] at (-0.18,1.35) {\(n\)};
\fill (0,0) circle (2pt);
\node[below right=1pt] at (0,0) {\(p\)};
\node[orange!85!black,font=\small] at (2.55,1.05) {沿这个方向朝 \(n\) 弯};
\node[blue!70!black,font=\small] at (-2.35,-1.35) {沿正交方向朝反侧弯};
\end{tikzpicture}
\end{center}

\emph{在同一点朝不同方向切下去，得到的曲线弯法不同，一个数装不下。}

\subsection{一个数说不清曲面怎样弯}

圆柱面把矛盾摆得最清楚：它在 \(\R^{3}\) 里明显弯着，而它上面所有的内部测量——长度、角度、面积——与平面完全一致。若坚持用一个数刻画弯曲，这个数该说圆柱弯还是平？

两个答案各有各的道理，因为它们回答的是两个问题。内蕴的弯曲是住在曲面上的居民能测到的；外在的弯曲是从环境空间看到的。圆柱面外在弯、内蕴平。球面两样都弯——这也是为什么纸能卷成圆筒却包不住皮球，包装纸在球面上必定起褶。

\textbf{第一基本形式记录内蕴，第二基本形式记录外在。} 本节余下的部分就是把这两句话写成公式。

\subsection{第一基本形式量的是曲面内部能测到的东西}

设曲面 \(S\subset\R^{3}\) 由参数化 \(X(u,v)\) 给出，切平面由 \(X_{u},X_{v}\) 张成。把环境的欧氏内积限制到切平面，在这组基下写成矩阵：

\[
\mathrm{I}=\begin{pmatrix}E&F\\F&G\end{pmatrix},
\qquad
E=\ip{X_{u}}{X_{u}},\quad
F=\ip{X_{u}}{X_{v}},\quad
G=\ip{X_{v}}{X_{v}}.
\]

这就是上一单元说的诱导度量，只不过换了个古典的名字和记号。曲面上一条参数曲线的长度是 \(\int\sqrt{E\,du^{2}+2F\,du\,dv+G\,dv^{2}}\)，两个切向量的夹角由 \(\mathrm{I}\) 定义的内积算出，面积元是 \(\sqrt{EG-F^{2}}\,du\,dv\)。凡是能用尺和量角器在曲面内部完成的测量，答案全在 \(\mathrm{I}\) 里。

球面 \(X(\theta,\phi)=(R\sin\theta\cos\phi,\,R\sin\theta\sin\phi,\,R\cos\theta)\) 的第一基本形式是

\[
\mathrm{I}=\begin{pmatrix}R^{2}&0\\0&R^{2}\sin^{2}\theta\end{pmatrix}.
\]

\(F=0\) 说明经线与纬线处处正交；\(\sin^{2}\theta\) 说明纬向的尺度随着靠近两极而收缩。等距圆柱投影在极地把东西方向拉得离谱，度量层面的原因就写在这个 \(\sin^{2}\theta\) 上：投影把它硬按成常数。

内蕴几何要问的全部问题可以压成一句：哪些量只依赖 \(\mathrm{I}\)。长度、角度、面积、测地线都只依赖它；``这张曲面在 \(\R^{3}\) 里摆成什么姿势''不依赖它。

\subsection{第二基本形式量的是曲面在外部空间里怎样弯}

外在的弯曲藏在法向的摆动里。取单位法向

\[
n=\frac{X_{u}\times X_{v}}{\norm{X_{u}\times X_{v}}},
\]

平面上它恒定不动，球面上它随点转个不停。沿切方向 \(w\) 移动时法向的变化率定义形状算子 \(S(w)=-dn(w)\)，负号只是约定。\(S\) 是切平面到自身的线性映射，它在基 \(X_{u},X_{v}\) 下的矩阵配上第一基本形式，给出第二基本形式

\[
\mathrm{II}=\begin{pmatrix}e&f\\f&g\end{pmatrix},
\qquad
e=\ip{X_{uu}}{n},\quad
f=\ip{X_{uv}}{n},\quad
g=\ip{X_{vv}}{n}.
\]

直觉不难说：\(X_{uu}\) 是参数曲线的加速度，\(e\) 取它在法向上的分量，也就是曲面沿 \(u\) 方向脱离切平面的速度。切向分量被丢掉了，因为那部分只反映参数化的走法，与形状无关——这跟上一节要求单位速度是同一个动作。平面的 \(\mathrm{II}\) 恒为零；半径 \(R\) 的球面有 \(\mathrm{II}=-\mathrm{I}/R\)，法向变化率处处一样，与半径成反比。

给定切方向 \(w\)，法曲率是两个形式的商

\[
\kappa_{n}(w)=\frac{\mathrm{II}(w,w)}{\mathrm{I}(w,w)},
\]

它就是把曲面沿含 \(n\) 与 \(w\) 的平面切开，所得截线的曲率（带符号）。分母的作用是抵消 \(w\) 的长度，于是 \(\kappa_{n}\) 只依赖方向。开头那张马鞍面的图画的正是它：同一点上，橙色方向的截线朝 \(n\) 一侧弯，蓝色方向的截线朝另一侧弯，\(\kappa_{n}\) 一正一负。方向不同答案不同，所以法曲率不是曲面的一个标量属性。

\subsection{两个主曲率，两种平均}

形状算子是切平面上的自伴算子，可以对角化，特征值 \(\kappa_{1}\ge\kappa_{2}\) 就是法曲率的最大值与最小值，称主曲率；对应的特征方向互相正交，称主方向。曲面在一点的外在弯曲，被两个数和两个方向说完了。

两种组合最常用：

\[
K=\kappa_{1}\kappa_{2}=\frac{eg-f^{2}}{EG-F^{2}}\quad(\text{Gauss 曲率}),
\qquad
H=\frac{\kappa_{1}+\kappa_{2}}{2}\quad(\text{平均曲率}).
\]

一个是形状算子的行列式，一个是它的迹的一半。椭球面上两个主曲率同号，\(K>0\)；马鞍面上异号，\(K<0\)；圆柱面沿母线方向的主曲率为零，于是 \(K=0\)，而 \(H=\kappa_{1}/2\neq0\)。

平均曲率有它自己的用武之地。\(H\equiv0\) 的曲面叫极小曲面，肥皂膜就是实例——两个方向的弯曲恰好抵消，面积在局部达到最小。平均曲率流 \(\partial_{t}X=Hn\) 让曲面沿法向以 \(H\) 的速率收缩，是几何分析里的基本方程，也是图像处理中``曲率驱动的平滑''的来源。

\subsection{内蕴与外在的分工，是这一单元的主线}

把两张清单摆在一起，界线就清楚了。只依赖 \(\mathrm{I}\) 的量是内蕴的：长度、角度、面积、测地线。依赖嵌入方式的量是外在的：法向、形状算子、主曲率、主方向、平均曲率。

分界线的意义由等距给出。两个曲面局部等距，指的是存在一个双射把一方的 \(\mathrm{I}\) 搬成另一方的 \(\mathrm{I}\)；此时住在两张曲面上的居民做任何内部测量都得到相同结果，他们无法通过测量区分自己在哪一张上。平面与圆柱面正是这样一对——展开圆柱的过程一寸都没拉伸，\(\mathrm{I}\) 原样保留。但它们的 \(\mathrm{II}\) 天差地别：平面的恒为零，圆柱的不是。同一套内蕴几何，可以有很多种外在实现。

这个分工也是数据视角里``低维结构''的几何原型。流形假设关心的是数据流形自己的度量与曲率，不是它恰好被嵌进了哪个高维空间；换一组特征、换一个可逆的线性变换，嵌入方式变了，内蕴结构可以一点没变。Isomap 估计的测地距离是内蕴量，Fisher 度量诱导的统计几何同样是内蕴量——同一个分布族换一种参数化，参数空间看起来完全不同，统计结构却是同一个。这些是类比，不是定理；但它们提示了一个值得反复问的问题：你手上的那个量，换个嵌入还在不在。

\subsection{Gauss 曲率的定义里明明用到了第二基本形式}

最后留一个悬案。按定义，Gauss 曲率 \(K=(eg-f^{2})/(EG-F^{2})\) 的分子来自第二基本形式，而第二基本形式是外在的——它需要法向，法向需要环境空间。照这个逻辑，\(K\) 应该是外在量，住在曲面上的居民没法测。

事实相反。下一节\hyperref[sec:12-Real-Analysis-Support/11-Curve-and-Surface-Geometry/03]{``从 Gauss 曲率到统计与数据流形''}要讲的绝妙定理说，\(K\) 只依赖 \(\mathrm{I}\)：一个看起来关于``怎样嵌入外部空间''的量，其实纯粹内蕴。这个结论会顺手解释一件日常的事——为什么球面永远画不成一张不失真的平面地图。
